% Experimental setup for the v0.6.0 sweep. Design facts only; every result
% value is \pending{} (defined in paper.tex).

\section{Experimental Setup}\label{sec:setup_exp}

\subsection{Data and provenance}\label{subsec:datasets}

We use two published K562 Perturb-seq screens in the scPerturb
repackaging~\citep{peidli2024scperturb} (Zenodo record $13350497$):
\textsc{Adamson 2016}~\citep{adamson2016perturbseq}, a CRISPR interference
screen, from its three 10X runs (pilot, 10X005, 10X010) concatenated onto the
intersection of their gene vocabularies; and
\textsc{Norman 2019}~\citep{norman2019perturbseq}, a CRISPR activation screen
with single and double perturbations. All four files are fetched
programmatically; each file's SHA-256 digest is pinned in
\texttt{src/perturb\_eval/data/download.py}, the digests were cross-checked
against the size and MD5 published in the Zenodo record, and a missing pin or
a digest mismatch refuses the fetch. Perturbation labels follow a recorded
label contract. Adamson raw construct labels are parsed structurally: the
plasmid suffix and \texttt{\_only} are stripped; a one-component construct
becomes a gene task (plasmids for the same gene are pooled and the guide
count per gene recorded); multi-gene constructs are excluded; negative-control
constructs are controls; cells with a missing perturbation annotation are
excluded and counted; and labels whose target cannot be corroborated in the
data are excluded with the reason. Norman symbols renamed since publication
are joined to the dataset vocabulary on their Ensembl gene IDs, and labels
that resolve under neither name are excluded with the reason. The held-out
task plan is drawn deterministically: Adamson targets are binned into $3$
quantile bins of mean $|\Delta\log1\mathrm{p}|$ on the target gene and $7$
are drawn per bin from the structurally eligible single-gene constructs
($21$ tasks); Norman contributes $15$ singletons and $5$ doublets ($20$
tasks), stratified by CRC32-derived strata, for $41$ tasks in total, all with
\texttt{seed=2026}. The sampler asserts that every stratum is filled exactly
and the run is refused otherwise. Every excluded label, every digest, the
resolved task plan with its strata and eligible-pool sizes, and the per-label
cell cap (\texttt{max\_cells\_per\_pert}) are written to the run's
provenance record (\S\ref{sec:reproducibility}).

\paragraph{On-target signal.} Each backbone receives the signed
log-fold-change of the perturbation's target gene(s) as its on-target
feature; for a doublet it is the mean over the two target genes. In Adamson
(CRISPR interference) the on-target signal is a decrease; in Norman (CRISPR
activation) it is an increase. Because the feature is signed, the direction
is learned from the data rather than assumed.

\subsection{Features and evaluation}\label{subsec:eval}

\paragraph{Features.} For each held-out task, the top-$2000$ highly variable
genes are selected on that task's \emph{training} cells only, with the target
gene(s) force-included. Models for different tasks therefore use different
feature spaces, and every cross-task median reported here is a median over
models with differing inputs.

\paragraph{Evaluation genes.} The error metric is the mean squared deviation
(MSD) between predicted and observed mean log-fold-change on the $20$ genes
with the largest absolute mean difference between the held-out
perturbation's cells and control cells --- the CPA/GEARS
convention~\citep{lotfollahi2023cpa,roohani2024gears}. These genes are chosen
using the held-out cells.

\subsection{Backbone pool}\label{subsec:backbones}

The Architect picks from three backbones: \textsc{linear} (ridge regression
per gene, NumPy); \textsc{mlp} (a two-layer MLP, PyTorch); and a
$2.1$\,M-parameter gene-token transformer trained from scratch (code
identifier \texttt{scgpt\_small}; not a pretrained foundation model).

\subsection{Trainer sweep and the oracle quantity}\label{subsec:oracle}

For every task the trainer sweep fits all $54$ configurations
$\{\textsc{linear}, \textsc{mlp}, \texttt{scgpt\_small}\} \times
N\in\{3,5\} \times R\in\{1,2,3\} \times 3$ seeds, with the seed passed to the
trainer. The quantity gated in H1/H2 is, per task, the minimum MSD over the
$54$ configurations, then its median over tasks. The minimum is selected on
the same evaluation genes and cells it is scored on, so it is an
\emph{oracle}: an attainable upper bound on what the configuration space can
reach, not an estimate of held-out performance. There is no nested
validation/test split.

\subsection{Agentic lifecycle}\label{subsec:lifecycle}

Five agents --- DataCurator, Literature, Architect, Trainer, Validator ---
produce Pydantic-validated proposals over the following configuration space:

\begin{itemize}
  \item \textbf{DataCurator}: \texttt{hvg\_method}$\in$\{seurat, scanpy\},
  \texttt{hvg\_count}$\in$\{500, 1000, 2000, 5000\},
  \texttt{qc\_mito\_max}, \texttt{split\_strategy}$\in$\{per\_pert\_holdout, unseen\_gene\},
  \texttt{batch\_correction}$\in$\{none, combat, harmony\}.
  \item \textbf{Literature}: BioGPT mechanism retrieval plus PubMed
  E-utilities and STRING-DB neighbour queries; output is a structured pathway
  prior and PPI neighbours for the Architect.
  \item \textbf{Architect}: backbone $\in$\{linear, mlp, scgpt\_small\},
  \texttt{n\_agents}$\in$[2,8], \texttt{n\_rounds}$\in$[1,5],
  \texttt{hvg\_count}, \texttt{learning\_rate}$>0$,
  \texttt{ridge\_lambda}$\ge 0$, \texttt{epochs}$\in$[1,500].
  \item \textbf{Trainer}: \texttt{lr}, \texttt{epochs},
  \texttt{ridge\_lambda}.
  \item \textbf{Validator}: dynamic threshold $\in$[0.02, 0.3] derived
  from the task's probe signature; returns a
  \texttt{StructuredCritique} with \texttt{which\_genes\_failed} and
  \texttt{suggested\_next\_config\_delta}, which the Architect applies in the
  next round.
\end{itemize}

\paragraph{LLM condition.} The agents are served by a rotating pool of
OpenRouter models with failover between pool members. The pool roster is
recorded in provenance, and every lifecycle step records the model that
served it (\texttt{model\_id}) and is labelled \texttt{source} $\in$
\{\texttt{llm}, \texttt{fallback}\}, where \texttt{fallback} marks the
rule-based default the pool returns after an LLM failure. A run with any
fallback step is invalid by construction: it is marked failed and the
analyser refuses to summarise it. Agent-choice statistics are computed over
LLM-sourced steps only. Each task runs the lifecycle with three seeds; the
seed is passed to the trainer and into the LLM cache key. The per-task seed
medians used for H4 and H5 depend on this: in the v0.5.0 artifacts the seed
did not reach the lifecycle run, and the three seeds were identical runs.

\paragraph{What the lifecycle trace records.} Each step records its role,
round, proposal, LLM confidence, \texttt{source} and \texttt{model\_id}. The
five roles do not critique each other, and no round has a winner. The trace
therefore has no critique matrix and no winner sequence, so CSD, WFR and the
four-component TDI of \S\ref{sec:metrics} cannot be computed on this testbed.
H4 and H5 use only the confidence-derived quantities
(\S\ref{sec:tdi-real}).

\subsection{Compute}\label{subsec:compute}

Preflight, trainer sweep and lifecycle sweep run in one process on a single
Modal A100-40GB container with an in-loop budget cap of \$28; the cap is a
limit, not the spend. Every run is appended atomically to
\texttt{trainer\_runs.jsonl} or \texttt{lifecycle\_runs.jsonl}, whose
record~0 is the run's provenance record. Actual spend and GPU-hours:
\pending{actual cost (USD) and GPU-hours from provenance.json}.
